\section{Capa de asociacion temporal basada en ByteTrack}
\label{sec:bytetrack-layer}

El proceso de detección permite identificar que en un frame aparece un elemento de una determinada clase, pero cada frame se procesa de forma independiente. Para que el sistema pueda reconocer a un mismo jugador a lo largo del vídeo, no basta con detectar jugadores, porteros, árbitros y balón en cada imagen: es necesario asociar las detecciones de frames consecutivos y mantener una identidad temporal estable. Este proceso se conoce como \textit{tracking multiobjeto}.

En nuestro caso, el tracking es una fase crítica porque gran parte de la información posterior depende del identificador asignado a cada elemento. Si un jugador recibe un ID distinto cada vez que aparece, no es posible construir su trayectoria, asociarle una posición, vincularlo a una alineación o utilizarlo correctamente en la generación de eventos. Por ello, el objetivo no es únicamente dibujar bounding boxes sobre el vídeo, sino mantener una representación persistente del partido. Idealmente, la salida final debe estar limitada a los elementos reales del juego: 20 jugadores de campo, 2 porteros, 3 árbitros y 1 balón.

Esta fase ha sido una de las más complejas del proyecto. En fútbol aparecen muchos problemas propios del dominio: oclusiones entre jugadores, cruces continuos, jugadores del mismo equipo con apariencia muy parecida, cambios de escala por el zoom, movimientos de cámara, detecciones perdidas durante varios frames y elementos que salen temporalmente del plano. Además, intentar solucionar uno de estos problemas de forma aislada puede empeorar otro. Por ejemplo, permitir reasociaciones muy amplias ayuda a recuperar jugadores tras una oclusión, pero también aumenta el riesgo de intercambiar IDs entre jugadores cercanos. Por ello, el sistema se ha construido como una combinación de varias restricciones complementarias.

La fase de asociacion temporal se ha dividido en dos capas. La primera capa, tratada en esta seccion, es una adaptacion de ByteTrack orientada a producir asociaciones robustas entre detecciones consecutivas. La segunda capa, descrita en la \hyperref[sec:tracking-canonico]{Seccion~\ref*{sec:tracking-canonico}}, intenta coser esas asociaciones de ByteTrack para estabilizar las 26 identidades finales del partido.

En la \hyperref[sec:tracking-tecnologias]{Seccion~\ref*{sec:tracking-tecnologias}} se introdujo ByteTrack desde un punto de vista general. En este apartado se describe su integracion concreta como primera capa de tracking del pipeline visual y las modificaciones realizadas sobre el algoritmo clásico.

\subsection{Base de la fase}

Sobre la estructura original de ByteTrack se han añadido señales propias del proyecto: clase reajustada, equipo estimado y posicion proyectada sobre el campo.

La fase consume las detecciones enriquecidas de identificacion y produce una salida intermedia de detecciones asociadas con \textit{raw tracker id}. Esta salida no es todavia la identidad final del jugador, sino una asociacion temporal que alimentará la siguiente capa.

\subsection{Asociacion multi-etapa en cada frame}

Manteniendo la logica general explicada en la \hyperref[sec:tracking-tecnologias]{Figura~\ref*{sec:tracking-tecnologias}}, la implementacion actual ejecuta seis subfases de matching en vez de las 3 descritas en \hyperref[sec:tracking-tecnologias]{2.4}:
Manteniendo la logica general de ByteTrack descrita en la \hyperref[sec:tracking-tecnologias]{Sección~\ref*{sec:tracking-tecnologias}} (asociacion progresiva de tracks activos, perdidos y no confirmados), nuestra adaptación desdobla el matching en seis subfases, ampliando el esquema base presentado en 2.4.
\begin{itemize}
    \item \texttt{high\_iou}
    \item \texttt{low\_iou}
    \item \texttt{unconfirmed\_iou}
    \item \texttt{high\_bbox}
    \item \texttt{low\_bbox}
    \item \texttt{unconfirmed\_bbox}
\end{itemize}

Mientras que las primeras tres subfases usan el solape entre cajas como función de coste, las últimas 3 subfases, utilizan la distancias entre los centros de las \textit{bounding box}, y actuan como rescate cuando el solape no es suficiente, reduciendo fragmentaciones en escenas con zoom, paneos y discontinuidades de deteccion.

\subsection{Restricciones por clase y equipo}

Otra mejora introducida sobre ByteTrack consiste en utilizar la clase del objeto y el equipo asignado como restricciones de asociación. Si un ID canónico corresponde a un jugador, no debería poder reasociarse con una detección de balón o de árbitro. Del mismo modo, si un jugador pertenece a un equipo, no debería heredar la identidad de un jugador del equipo contrario.

Por lo tanto, en nuestra adaptación del ByteTrack, se usan la clase de la detección y su equipo asociado como restricciones de asociación duras. Si la clase del track y la clase de la deteccion candidata no coinciden, la asociacion se descarta automáticamente.

Del mismo modo, para \texttt{player}, cuando hay informacion de equipo tanto en el track como en la detección, se aplica también una restricción dura: si no coinciden, la asociacion se descarta.

Estas restricciones reducen los intercambios entre elementos de distinta clase y entre jugadores de equipos diferentes. Sin embargo, todavía puede haber intercambios entre jugadores del mismo equipo, especialmente cuando se cruzan o aparecen muy juntos en la imagen. Para reducir este problema fue necesario incorporar información geométrica del campo.

\subsection{Uso de homografia y restricción espacial en metros}

En nuestra versión de ByteTrack adaptada también se aplica un gate espacial en campo para las clases \texttt{player} y \texttt{goalkeeper}. Si la distancia proyectada entre el track y la detección candidata supera el límite configurado (8 metros), la asociación se descarta. En cambio, si está dentro del rango permitido, la pareja se mantiene como factible y el coste se decide por el criterio de matching de la subfase (IoU o distancia de centros en las fases de rescate).

Durante el desarrollo se evaluó usar la posición de las detecciones proyectada al campo como coste principal de la asociación. En la practica, esta opción degradaba el rendimiento ya que esta señal presentaba mucha vibracion frame a frame debido a pequeños cambios en la homografía entre frames.  Por este motivo había variaciones bruscas en las trayectorias estimadas de las identidades. Esto afectaba especialmente a la prediccion del estado con el filtro de Kalman y acababa propagando errores al proceso de matching. Por este motivo, en la version final la posicion en campo se utiliza como compuerta geometrica de factibilidad (gate duro) y no como termino principal de coste.

\subsection{Control de creacion de nuevos tracks}

Otro problema relevante es la creacion excesiva de nuevos tracks. Esto puede ocurrir cuando el detector produce bounding boxes duplicados o cuando ByteTrack pierde temporalmente una identidad. Para reducir este efecto, se añadieron reglas de supresión por solape de tracks nuevos:
\begin{itemize}
    \item contra tracks activos ya consolidados,
    \item contra tracks tentativos no confirmados,
    \item y entre candidatos nuevos del mismo frame.
\end{itemize}

Un candidato a track nuevo se descarta si solapa demasiado con un track activo, si solapa con un track tentativo no consolidado, o si solapa con otro candidato nuevo de mayor confianza. Ante cualquier solape entre dos candidatos nuevos, se conserva únicamente el de mayor confianza. Con esta regla se consigue reducir tracks duplicados sin alterar el seguimiento de tracks ya consolidados.

La configuración concreta de estos umbrales de solape se resume en la \hyperref[tab:supresion-nuevos-tracks-bytetrack]{Tabla~\ref*{tab:supresion-nuevos-tracks-bytetrack}}.

\begin{table}[htbp]
    \centering
    \scriptsize
    \begin{tabular}{c|c|p{7.5cm}}
        \textbf{Concepto} & \textbf{Valor} & \textbf{Funcion} \\
        \hline\hline
        Solape maximo con tracks activos & 0.2 & Evita crear un nuevo ID sobre un track ya consolidado. \\
        Solape maximo con tracks tentativos & 0.1 & Evita duplicados sobre tracks aun no consolidados. \\
        Solape entre candidatos nuevos & 0.0 & Conserva solo el candidato nuevo de mayor confianza si hay solape. \\
        \hline
    \end{tabular}
    \caption{Reglas de supresion en la creacion de nuevos tracks de ByteTrack.}
    \label{tab:supresion-nuevos-tracks-bytetrack}
\end{table}

\subsection{Resumen de la capa}

En resumen, esta fase no se limita a ejecutar ByteTrack sobre cajas de YOLO, sino que integra restricciones de clase, equipo, geometria en campo y control de creación de tracks para mejorar la asociacion temporal en video de futbol broadcast. La capa siguiente toma estas asociaciones y construye las identidades canonicas estables del partido.
